% from ega2-1-fr.tex
\oldpage[II]{175}
Mais avec les notations de (\hyperref[II.8.6.2-fr]{8.6.2}), si
$v:C_X\to X$ est le morphisme structural, il résulte de
(\hyperref[II.8.6.2.1-fr]{8.6.2.1}) que l'on a
$v^{-1}(X_f)=\operatorname{Spec}(S_f^{\geqslant})$. On a d'autre part
$\operatorname{Spec}(S_{(f^{\natural})}^{\natural})=D_+(f^{\natural})$, si
bien que (\hyperref[II.8.7.3.2-fr]{8.7.3.2}) définit un isomorphisme
$v^{-1}(X_f)\to D_+(f^{\natural})$. En outre, si $g\in S_e$ ($e>0$), le
diagramme
\[
  \xymatrix{
    v^{-1}(X_{fg})\ar[r]^{\sim}\ar[d] &
    D_+(f^{\natural}g^{\natural})\ar[d]\\
    v^{-1}(X_f)\ar[r]^{\sim} & D_+(f^{\natural})
  }
\]
est commutatif, comme il résulte aussitôt de la définition de l'isomorphisme
(\hyperref[II.8.2.7.3-fr]{8.2.7.3}). Enfin par définition $S_+$ est engendré
par les $f$ homogènes, donc il résulte de
(\hyperref[II.8.2.10-fr]{8.2.10, (iv)}) et de
(\hyperref[II.2.3.14-fr]{2.3.14}) que les $D_+(f^{\natural})$ forment un
recouvrement de $\operatorname{Proj}(S^{\natural})$ et les $v^{-1}(X_f)$ un
recouvrement de $C_X$, puisque les $X_f$ forment un recouvrement de $X$~;
cela achève donc de définir dans ce cas l'isomorphisme
(\hyperref[II.8.7.3.1-fr]{8.7.3.1}).

Pour démontrer (\hyperref[II.8.7.3-fr]{8.7.3}) dans le cas général, il suffit
de voir que si $U$, $U'$ sont deux ouverts affines de $Y$ tels que
$U'\subset U$, d'anneaux $A$ et $A'$, et si on pose
$\mathcal{S}|U=\widetilde{S}$, $\mathcal{S}|U'=\widetilde{S'}$, le diagramme
\begin{equation}
\label{II.8.7.3.3-fr}
  \xymatrix{
    C_{U'}\ar[r]\ar[d] & \operatorname{Proj}(S'^{\natural})\ar[d]\\
    C_U\ar[r] & \operatorname{Proj}(S^{\natural})
  }
\tag{8.7.3.3}
\end{equation}
est commutatif. Mais $S'$ s'identifie canoniquement à $S\otimes_A A'$, donc
$S'^{\natural}$ à
\[
  S^{\natural}\otimes_S S'=S^{\natural}\otimes_A A'~;
\]
on a donc
$\operatorname{Proj}(S'^{\natural})=\operatorname{Proj}(S^{\natural})
\times_U U'$ (\hyperref[II.2.8.10-fr]{2.8.10})~; de même, si
$X=\operatorname{Proj}(S)$, $X'=\operatorname{Proj}(S')$, on a
$X'=X\times_U U'$ et
$\mathcal{S}_{X'}=\mathcal{S}_X\otimes_{\mathcal{O}_U}\mathcal{O}_{U'}$
(\hyperref[II.3.5.4-fr]{3.5.4}), autrement dit
$\mathcal{S}_{X'}=j^*(\mathcal{S}_X)$, où $j$ est la projection $X'\to X$.
On a par suite (\hyperref[II.1.5.2-fr]{1.5.2})
$C_{U'}=C_U\times_X X'=C_U\times_U U'$, et la commutativité de
(\hyperref[II.8.7.3.3-fr]{8.7.3.3}) est alors immédiate.
\end{proof}

\begin{remark}[8.7.4]
\label{II.8.7.4-fr}
\begin{enumerate}
  \item[\rm{(i)}] La fin du raisonnement de
  (\hyperref[II.8.7.3-fr]{8.7.3}) se généralise aussitôt de la façon
  suivante. Soient $g:Y'\to Y$ un morphisme,
  $\mathcal{S}'=g^*(\mathcal{S})$, $X'=\operatorname{Proj}(\mathcal{S}')$~;
  on a alors un diagramme commutatif
  \begin{equation}
  \label{II.8.7.4.1-fr}
    \xymatrix{
      C_{X'}\ar[r]\ar[d] &
      \operatorname{Proj}(\mathcal{S}'^{\natural})\ar[d]\\
      C_X\ar[r] & \operatorname{Proj}(\mathcal{S}^{\natural})
    }
  \tag{8.7.4.1}
  \end{equation}

  D'autre part, soit $\varphi:\mathcal{S}''\to\mathcal{S}$ un homomorphisme
  de $\mathcal{O}_Y$-Algèbres graduées, tel que si on pose
  $X''=\operatorname{Proj}(\mathcal{S}'')$,
  $u=\operatorname{Proj}(\varphi):X\to X''$ soit partout défini~; on a
  d'autre part

  